\documentclass[11pt]{article}

\usepackage[margin=1in]{geometry}
\usepackage{array}
\usepackage{booktabs}
\usepackage{xurl}
\usepackage{hyperref}
\usepackage{xcolor}
\usepackage{enumitem}
\usepackage{longtable}

\hypersetup{
  colorlinks=true,
  linkcolor=blue,
  urlcolor=blue,
  citecolor=blue
}

\setlist[itemize]{leftmargin=1.25em}
\setlist[enumerate]{leftmargin=1.5em}
\newcommand{\code}[1]{\texttt{\small #1}}
\newcommand{\filepath}[1]{\begingroup\urlstyle{tt}\small\url{#1}\endgroup}

\title{\textbf{Professor Q\&A Notes: Taxonomy-aware BiomeGPT, Prediction Results, and Batch Correction}}
\author{Yuanshu / Codex working report}
\date{May 20, 2026}

\begin{document}
\maketitle

\section*{How To Use This Document}

This is not the main technical report. It is a meeting-preparation document: likely
questions, direct answers, and the reasoning behind our choices. The main message is that
we have three connected pieces of work:

\begin{enumerate}
  \item A taxonomy-aware BiomeGPT model for microbiome representation learning.
  \item A Healthy-vs-Diseased prediction pipeline with external validation.
  \item A batch-effect diagnostic and preliminary adversarial correction pipeline.
\end{enumerate}

\section{One-Minute Summary}

\textbf{What did we build?} We reproduced a BiomeGPT-style microbiome transformer and
extended it with taxonomy-aware species embeddings. Each sample is represented by a
sample-level \code{<cls>} embedding, and each microbial species has a learned species
embedding informed by taxonomy.

\textbf{Why taxonomy-aware?} Microbial species are not independent IDs. Species from the
same genus, family, order, or phylum share biology. Adding rank-wise taxonomy embeddings
gives the model an inductive bias that should improve interpretability and transfer when
external cohorts have sparse or shifted species coverage.

\textbf{Why batch diagnostics?} External validation can fail because the model learns
study/cohort signals instead of robust disease biology. We reconstructed study/cohort
batch-proxy labels from public accession metadata and tested whether the pretrained
embeddings encode those labels.

\textbf{What did we find?} Batch signal is present, but broad batch labels are strongly
phenotype-confounded. Naive adversarial correction using broad labels can fail or even
increase post-hoc batch separability. Conservative-safe labels remain the best correction
target, but adversarial optimization needs careful scheduling.

\section{What Is The Biological Task?}

\textbf{Likely question: What biological problem are we solving?}

We are studying whether gut microbiome species-abundance profiles can be represented by a
foundation-model style transformer and used for Healthy-vs-Diseased prediction. A gut
microbiome sample is a sparse vector of microbial species abundances. The biological
challenge is that disease-associated microbial patterns are subtle, cohort-dependent, and
strongly affected by geography, sequencing protocol, study design, and phenotype mixture.

\textbf{Answer to give:}

\begin{quote}
We are not only training a binary classifier. We are trying to learn a transferable
microbiome representation: species-level embeddings and sample-level embeddings that
capture community structure, taxonomy, and disease-associated shifts across cohorts.
\end{quote}

\section{What Is BioGPT / BiomeGPT Doing?}

\textbf{Likely question: How does the model represent a microbiome sample?}

In our implementation, each sample is a sequence of species tokens. Each token has two
parts:

\begin{itemize}
  \item a species identity representation;
  \item an abundance-bin representation.
\end{itemize}

The model uses masked abundance reconstruction during pretraining. We randomly mask a
fraction of observed nonzero species abundances and ask the transformer to reconstruct
their abundance bins from the rest of the microbial community. This is analogous to masked
language modeling, except the tokens are microbial species and the prediction target is an
abundance bin rather than a word.

\textbf{Key output representations:}

\begin{itemize}
  \item \textbf{Species prompt}: learned species representation, analogous to a gene
  embedding in scGPT.
  \item \textbf{Sample prompt}: final \code{<cls>} embedding for a microbiome sample,
  analogous to a cell embedding in scGPT.
\end{itemize}

\section{What Did We Change: Taxonomy-aware BiomeGPT}

\textbf{Likely question: What is new compared with a plain species-ID model?}

The original species token can treat every species as an unrelated ID. We changed the
species representation to include taxonomy:

\begin{quote}
\code{Domain + Kingdom + Phylum + Class + Order + Family + Genus + Species}
\end{quote}

This means two species from the same genus share the genus component; species from the
same family share the family component; and so on.

\textbf{Why this is biologically reasonable:}

\begin{itemize}
  \item Microbial traits are partially structured by taxonomy.
  \item External datasets may contain related but not identical species.
  \item Taxonomy-aware embeddings make the species latent space more interpretable.
  \item The inductive bias may help transfer when species overlap across cohorts is
  imperfect.
\end{itemize}

\section{What Data Did We Use?}

\begin{center}
\begin{tabular}{lr}
\toprule
Data component & Count / role \\
\midrule
Phase-1 pretraining species vocabulary & 1,012 species \\
Phase-2 gut samples for pretraining / batch diagnostics & 13,332 samples \\
Fine-tuning training set & \code{\_prev3} gut data \\
External validation set & 927 samples \\
ExVal label balance & 612 Healthy / 315 Diseased \\
Conservative-safe batch-correction samples & 2,134 samples \\
\bottomrule
\end{tabular}
\end{center}

The ExVal set overlaps with 512 of the 513 \code{\_prev3} fine-tuning species, so external
evaluation is feasible after reindexing species columns to the training vocabulary.

\section{Pipeline Overview}

\begin{enumerate}
  \item \textbf{Phase 1 pretraining}: gut + non-gut samples, learn broad microbiome
  ecological structure.
  \item \textbf{Phase 2 pretraining}: gut-only adaptation, make the representation more
  domain-specific for stool/gut samples.
  \item \textbf{Embedding extraction}: extract species prompts and sample prompts.
  \item \textbf{Healthy-vs-Diseased fine-tuning}: add classifier on sample prompt and train
  on \code{\_prev3}.
  \item \textbf{External validation}: evaluate on ExVal using accuracy, F1, AUROC,
  macro-accuracy, macro-F1, and class-specific H/D accuracy.
  \item \textbf{Batch diagnostics}: test whether phase-2 sample prompts encode
  study/cohort labels.
  \item \textbf{Adversarial correction}: preliminary GRL-based correction on
  conservative-safe batch labels.
\end{enumerate}

\section{Important Parameters And What They Mean}

\begin{longtable}{p{0.25\textwidth}p{0.65\textwidth}}
\toprule
Parameter & Meaning \\
\midrule
\endhead
\code{bins = 32} &
Continuous relative abundances are converted into 32 rank-based abundance bins. This
reduces scale sensitivity across datasets and turns abundance prediction into a structured
masked reconstruction problem. \\
\code{mask\_ratio = 0.25} &
During pretraining, 25\% of nonzero observed species are masked and reconstructed. This
forces the model to learn community context rather than memorizing each input value. \\
\code{d\_model = 512} &
Hidden embedding dimension for species/sample representations. Larger values can capture
more structure but increase training cost and overfitting risk. \\
\code{nhead = 8} &
Number of transformer attention heads. Multiple heads allow the model to learn different
interaction patterns among species. \\
\code{num\_layers = 8} &
Number of transformer blocks. More layers increase capacity but also training instability
and overfitting risk. \\
\code{lr} &
Learning rate. Pretraining used a larger learning rate than fine-tuning/correction because
fine-tuning starts from an already learned checkpoint. \\
\code{threshold} &
Probability cutoff for calling a sample Diseased in H/D prediction. This strongly affects
Healthy accuracy vs Diseased accuracy. \\
\code{macro-F1} &
Average of class-wise F1. Important because Healthy and Diseased counts are imbalanced. \\
\code{AUROC} &
Threshold-independent ranking metric. It can be good even if a chosen threshold gives poor
Healthy/Diseased balance. \\
\code{DAB weight} &
Weight on the adversarial batch-discriminator loss. Larger is not automatically better. \\
\code{GRL lambda} &
Strength of gradient reversal. It controls how strongly the encoder is pushed to hide
batch information. \\
\code{min\_batch\_size = 10} &
Minimum samples required for a batch label to be used in a probe/correction trial. This
avoids unstable tiny classes. \\
\end{longtable}

\section{Why Did We Need Batch Labels?}

\textbf{Likely question: Did the original metadata contain batch labels?}

No. The original phase-2 gut metadata did not contain explicit sequencing batch, study ID,
cohort ID, sequencing center, platform, country, library prep, or run-date columns. It
contained phenotype and sample IDs. Therefore, we reconstructed a study/cohort-level
batch-proxy annotation.

\textbf{How we reconstructed labels:}

\begin{itemize}
  \item Public accession IDs such as \code{SAMEA}, \code{SAMN}, \code{SAMD},
  \code{SRR}, and \code{EGAR} were queried against ENA, NCBI SRA/BioSample, and EGA.
  \item Local IDs without public accessions were assigned prefix-derived labels only when
  the pattern was structured and interpretable.
  \item Every label received confidence: high, medium, or low.
  \item Every batch label was checked for phenotype confounding.
\end{itemize}

\begin{center}
\begin{tabular}{lrl}
\toprule
Confidence & Samples & Interpretation \\
\midrule
High & 9,389 & Exact public accession lookup returned study/cohort metadata \\
Medium & 3,418 & Structured local prefix-derived cohort label \\
Low & 525 & Unresolved or ambiguous local sample ID pattern \\
\bottomrule
\end{tabular}
\end{center}

\section{What Is Batch-Phenotype Confounding?}

\textbf{Likely question: If the batch label is real, why not correct it directly?}

A real study label is not automatically safe. Many microbiome studies are phenotype
specific. For example, one study may contain only Healthy samples, while another study may
contain mostly Parkinson's disease or colorectal cancer samples. If we remove study signal
naively, we may also remove true disease biology.

\textbf{Rule used:}

\begin{quote}
\code{phenotype\_confounding\_warning=True} if the top phenotype fraction in a batch is
\code{>= 0.80} or if the batch has only one phenotype.
\end{quote}

\textbf{Observed confounding:}

\begin{center}
\begin{tabular}{lrrrr}
\toprule
Label panel & Cramer's V & NMI & Median top phenotype frac. & Batches top frac. $\geq$ 0.8 \\
\midrule
High + medium & 0.802 & 0.519 & 1.000 & 79 / 105 \\
High-only & 0.907 & 0.489 & 1.000 & 66 / 83 \\
Conservative-safe & 0.683 & 0.457 & 0.549 & 0 / 17 \\
\bottomrule
\end{tabular}
\end{center}

\textbf{Answer to give:}

\begin{quote}
We can use broad labels for diagnostics, but final correction should use only labels that
are high-confidence and not strongly phenotype-confounded.
\end{quote}

\section{What Is A Batch Probe?}

\textbf{Likely question: What does the batch probe actually measure?}

A batch probe is a shallow classifier trained on frozen sample embeddings. The model
itself is not changed. If a logistic regression classifier can predict batch labels from
the \code{<cls>} embedding, then the embedding contains batch-associated information.

\begin{center}
\begin{tabular}{lrrrr}
\toprule
Probe & Samples & Classes & Balanced acc. & Macro-F1 \\
\midrule
High + medium batch & 12,758 & 105 & 0.397 & 0.365 \\
High-only batch & 9,340 & 83 & 0.386 & 0.355 \\
Conservative-safe batch & 2,134 & 17 & 0.503 & 0.499 \\
Healthy vs Diseased probe & 13,332 & 2 & 0.801 & 0.796 \\
\bottomrule
\end{tabular}
\end{center}

\textbf{Important distinction:} the H/D probe is not the final ExVal classifier. It is a
diagnostic test showing that frozen pretrained embeddings contain H/D-separable signal in
the phase-2 data.

\section{How Did We Implement Batch Correction?}

We adapted scGPT-style domain-adversarial training:

\begin{itemize}
  \item Attach a small batch discriminator to the \code{<cls>} sample embedding.
  \item Use a Gradient Reversal Layer (GRL).
  \item Keep the original masked abundance reconstruction loss.
  \item Add a batch-discriminator loss with small weight.
\end{itemize}

The light correction objective was:

\begin{quote}
\code{total loss = reconstruction loss + 0.1 * batch cross-entropy loss}
\end{quote}

Because the batch branch uses GRL, the discriminator learns to predict batch, while the
encoder receives reversed gradients and is pushed to make \code{<cls>} less
batch-predictive.

\section{What Did Batch Correction Show?}

\begin{center}
\small
\begin{tabular}{p{0.44\textwidth}rrrr}
\toprule
Run & Batch F1 before & Batch F1 after & H/D F1 before & H/D F1 after \\
\midrule
Conservative-safe, batch size 8, DAB 0.1, 3 epochs & 0.499 & 0.476 & 0.796 & 0.789 \\
High + medium, batch size 32, DAB 0.1, 3 epochs & 0.365 & 0.383 & 0.796 & 0.804 \\
High-only, batch size 32, DAB 0.1, 3 epochs & 0.355 & 0.373 & 0.796 & 0.802 \\
Conservative-safe, batch size 32, DAB 0.1, 3 epochs & 0.499 & 0.511 & 0.796 & 0.794 \\
Conservative-safe, DAB 0.5, 5 epochs & 0.499 & 0.505 & 0.796 & 0.786 \\
\bottomrule
\end{tabular}
\end{center}

\textbf{Interpretation:} the first conservative-safe light run moved in the desired
direction, but broader labels, larger batch size, and stronger DAB did not. This suggests
the issue is not simply ``train more.'' The adversarial game is sensitive to label quality,
batch size, and schedule.

\section{Improved Experiment: Evaluate The Pretraining Objective Directly}

\textbf{Likely question: Since BiomeGPT is a foundation model, should we evaluate batch
correction only by H/D prediction?}

No. H/D prediction is a downstream task. The core BiomeGPT pretraining task is masked
abundance prediction. Therefore, batch correction should also be judged by whether the
model preserves or improves masked abundance reconstruction, especially across
study/cohort labels.

I added a reconstruction diagnostic:

\begin{quote}
\filepath{dataset_v3/batch_reconstruction_diagnostics.py}
\end{quote}

It uses the same deterministic mask for each checkpoint, masks nonzero species positions,
and evaluates mean squared error on the true abundance bins. It also computes per-batch
MSE, batch MSE coefficient of variation, and batch MSE gap.

\begin{center}
\small
\begin{tabular}{p{0.34\textwidth}rrrr}
\toprule
Checkpoint & Overall MSE & Safe-panel MSE & Safe-panel CV & Safe-panel gap \\
\midrule
Original phase-2 & 34.857 & 37.808 & 0.095 & 13.140 \\
Conservative-safe, batch size 8 & 33.745 & 35.724 & 0.093 & 12.961 \\
Conservative-safe, batch size 32 & 34.088 & 36.252 & 0.096 & 13.324 \\
High + medium, batch size 32 & 33.049 & 35.770 & 0.096 & 12.736 \\
High-only, batch size 32 & 33.373 & 36.061 & 0.098 & 13.453 \\
Conservative-safe, DAB 0.5 & 33.853 & 35.256 & 0.092 & 12.554 \\
\bottomrule
\end{tabular}
\end{center}

\textbf{Interpretation:} reconstruction MSE generally improved after continued
correction training, so these trials did not destroy the masked-abundance objective.
However, lower reconstruction MSE alone is not enough: some runs improved reconstruction
while the batch probe got worse. Therefore, model selection should require three
conditions:

\begin{enumerate}
  \item batch probe decreases;
  \item masked abundance reconstruction is preserved or improves;
  \item H/D probe or downstream ExVal does not collapse.
\end{enumerate}

\textbf{Professor-facing wording:}

\begin{quote}
Because BiomeGPT is pretrained by masked abundance prediction, batch correction should not
be evaluated only through disease classification. The primary foundation-model check is
whether the corrected model remains good at masked abundance reconstruction, ideally with
more balanced reconstruction error across batches. Disease prediction should be a
downstream sanity check.
\end{quote}

\section{Why Was Prediction Around 0.5 Before, But Now Some Numbers Are 0.7--0.8?}

\textbf{This is the most important clarification. These numbers are not all the same
metric.}

\begin{center}
\small
\begin{tabular}{lp{0.45\textwidth}p{0.22\textwidth}}
\toprule
Number & What it means & Comparable to ExVal macro-F1? \\
\midrule
0.559 & ExVal macro-F1 from selected full run at threshold 0.6 & Yes \\
0.592 & Fixed-threshold ExVal macro-F1 at threshold 0.8 & Yes \\
0.609 & Best value in threshold sensitivity diagnostic at threshold 0.9 & Partly; diagnostic threshold sweep \\
0.769--0.783 & ExVal AUROC & No; threshold-independent ranking metric \\
0.796 & Frozen H/D probe macro-F1 on phase-2 embeddings & No; diagnostic probe, not ExVal prediction \\
0.96 & Internal training-side calibration macro-F1 during selection & No; internal split, not external validation \\
\bottomrule
\end{tabular}
\end{center}

\textbf{Answer to give:}

\begin{quote}
The earlier 0.5--0.6 values are external validation macro-F1, which is the hardest and
most relevant prediction metric. The 0.77--0.78 values are AUROC, not macro-F1. The 0.79
value is a frozen embedding H/D probe on phase-2 data, not the final external prediction.
So these numbers are not contradictory; they measure different questions.
\end{quote}

\textbf{Why ExVal macro-F1 is lower than the internal calibration score:}

\begin{itemize}
  \item ExVal is a different external cohort with study/domain shift.
  \item The classifier tends to over-call Diseased at lower thresholds, causing low Healthy
  accuracy.
  \item Macro-F1 penalizes imbalance between Healthy and Diseased performance.
  \item Batch/study effects may make the learned disease signal less transferable.
\end{itemize}

\section{How To Explain The ExVal Result}

The selected run had:

\begin{center}
\begin{tabular}{lr}
\toprule
Metric & Value \\
\midrule
ExVal accuracy & 0.563 \\
ExVal F1 & 0.602 \\
ExVal AUROC & 0.783 \\
ExVal macro-F1 & 0.559 \\
Healthy accuracy & 0.353 \\
Diseased accuracy & 0.971 \\
Threshold & 0.6 \\
\bottomrule
\end{tabular}
\end{center}

The fixed threshold 0.8 run improved Healthy specificity:

\begin{center}
\begin{tabular}{lr}
\toprule
Metric & Value \\
\midrule
ExVal accuracy & 0.593 \\
ExVal AUROC & 0.770 \\
ExVal macro-F1 & 0.592 \\
Healthy accuracy & 0.407 \\
Diseased accuracy & 0.956 \\
\bottomrule
\end{tabular}
\end{center}

\textbf{Interpretation:} the model ranks samples reasonably well by disease probability
(AUROC around 0.77--0.78), but calibration/thresholding is poor under external shift. It
predicts Diseased too often, which gives high Diseased accuracy but low Healthy accuracy.
This is why macro-F1 is modest.

\section{What If Professor Says Their Macro-F1 Is Around 0.7?}

\textbf{Answer carefully:}

\begin{quote}
That may be true, but we need to confirm the exact evaluation protocol: split, external
cohort, threshold selection, class labels, species alignment, whether metrics are macro-F1
or AUROC, and whether threshold was tuned on external labels. Our current fully external
macro-F1 is around 0.56--0.61 depending on threshold. We should not compare that directly
to internal validation, AUROC, or an externally tuned threshold.
\end{quote}

\textbf{What to ask:}

\begin{itemize}
  \item Was the 0.7 value macro-F1, weighted F1, accuracy, or AUROC?
  \item Was it measured on ExVal or on an internal split?
  \item Was threshold selected on the training side only, or tuned using ExVal labels?
  \item Were Healthy/Diseased labels constructed in the same way?
  \item Were species columns aligned to the same vocabulary?
  \item Did they use the same phase-2 checkpoint and the same \code{\_prev3} fine-tuning
  set?
\end{itemize}

\section{What Should We Do If Batch And Phenotype Are Confounded?}

\textbf{Do not blindly remove all batch signal.} If batch and phenotype are nearly the
same variable, adversarial correction can erase disease biology.

\textbf{Recommended approach:}

\begin{enumerate}
  \item Use broad labels only for diagnostics.
  \item Use high-confidence, non-confounded labels for correction.
  \item Keep post-correction H/D probes to check whether biological signal survives.
  \item Use warm-up schedules for DAB/GRL instead of fixed full adversarial pressure.
  \item Select correction checkpoints by a joint criterion: lower batch probe and
  preserved H/D probe.
  \item Only evaluate final value by downstream ExVal H/D metrics.
\end{enumerate}

\section{Why Broad Batch Correction Failed}

Our ablation suggests several possible mechanisms:

\begin{itemize}
  \item Broad batch labels are heterogeneous and phenotype-confounded.
  \item The discriminator-encoder adversarial game may be unstable.
  \item The encoder may learn a new representation that is still batch-separable to a fresh
  post-hoc probe.
  \item There was no warm-up schedule; adversarial pressure was applied immediately.
  \item Batch size changed adversarial dynamics; batch size 8 and 32 behaved differently.
\end{itemize}

\textbf{Meeting wording:}

\begin{quote}
The negative ablation is still informative. It shows that batch correction is not a
one-line improvement. Label quality and phenotype confounding dominate. We need a cautious
correction target and a scheduled adversarial optimization.
\end{quote}

\section{What Is The Best Next Experiment?}

\begin{itemize}
  \item Label panel: conservative-safe only.
  \item Batch size: 8 or 16.
  \item DAB weight: 0.02, 0.05, 0.1.
  \item GRL/DAB warm-up over first 30--50\% of epochs.
  \item Early stopping by lower safe-batch probe plus preserved H/D probe.
  \item Then fine-tune H/D from the best corrected checkpoint and evaluate ExVal.
\end{itemize}

\section{Short Answers To Likely Questions}

\begin{longtable}{p{0.34\textwidth}p{0.58\textwidth}}
\toprule
Question & Short answer \\
\midrule
\endhead
Is this exact sequencing batch? &
No. It is an externally enriched study/cohort batch-proxy annotation. Exact technical
batch would require original run date, kit, center, instrument, and processing metadata. \\
Why use batch proxies then? &
Because they are enough to test whether embeddings encode study/cohort signal, which is a
major source of external validation shift. \\
Why not correct all high-confidence labels? &
Because high-confidence study labels can still be disease-confounded. A real label can be
unsafe. \\
What does batch probe prove? &
It proves batch labels are predictable from frozen embeddings; it does not by itself prove
causality or solve correction. \\
Why did broad adversarial correction increase batch probe? &
Likely because labels are heterogeneous/confounded and the adversarial game can reorganize
batch information rather than remove it. \\
Is batch correction finished? &
No. We have a working diagnostic and proof-of-concept pipeline, but final correction needs
schedule tuning and ExVal validation. \\
Why are 0.7--0.8 numbers appearing? &
They are often AUROC or embedding-probe metrics, not ExVal macro-F1. The strict external
macro-F1 is currently around 0.56--0.61 depending on threshold. \\
What is the most defensible current claim? &
Taxonomy-aware BiomeGPT learns biologically meaningful sample embeddings, but external
prediction is limited by calibration/domain shift; batch correction is motivated but must
be conservative. \\
\bottomrule
\end{longtable}

\end{document}
